\section{Loop Invariant dan Verifikasi Perulangan}

Sebagian besar waktu eksekusi suatu algoritma dihabiskan di dalam struktur perulangan (\textit{loop}). Karena jumlah iterasi tidak selalu diketahui secara pasti sebelum eksekusi, diperlukan mekanisme formal untuk memverifikasi bahwa algoritma menghasilkan keluaran yang benar di setiap iterasi. Mekanisme tersebut adalah \logic{loop invariant}, yang menerapkan prinsip induksi matematika pada verifikasi program \cite{rosen2019}.

\subsection{Definisi Loop Invariant}

\begin{definisi}[Loop Invariant]
\logic{Loop invariant} adalah pernyataan (predikat) $I$ yang dikaitkan dengan struktur perulangan dan memenuhi tiga syarat:
\begin{enumerate}
    \item \textbf{Inisialisasi} (\textit{initialization}): $I$ bernilai benar sebelum iterasi pertama dimulai.
    \item \textbf{Pemeliharaan} (\textit{maintenance}): Jika $I$ bernilai benar sebelum suatu iterasi dan kondisi perulangan $B$ bernilai benar, maka $I$ tetap bernilai benar setelah iterasi tersebut dijalankan.
    \item \textbf{Terminasi} (\textit{termination}): Ketika perulangan berakhir ($\neg B$), invarian $I$ bersama dengan $\neg B$ mengimplikasikan postkondisi yang diinginkan.
\end{enumerate}
\end{definisi}

\subsection{Aturan Perulangan dalam Logika Hoare}

Untuk perulangan \texttt{while $B$ do $S$}:

\[ \frac{\{I \wedge B\}\ S\ \{I\}}{\{I\}\ \texttt{while } B \texttt{ do } S\ \{I \wedge \neg B\}} \]

Jika kita dapat menunjukkan bahwa badan perulangan $S$ menjaga invarian $I$ (yaitu, $\{I \wedge B\}\ S\ \{I\}$ valid), maka setelah perulangan berakhir, kita memperoleh $I \wedge \neg B$.

\subsection{Analogi dengan Induksi Matematika}

Hubungan antara loop invariant dan induksi matematika dapat dirinci secara berpasangan sebagai berikut:

\begin{center}
\renewcommand{\arraystretch}{1.3}
\begin{tabular}{|>{\raggedright\arraybackslash}p{5cm}|>{\raggedright\arraybackslash}p{5cm}|}
\hline
\textbf{Induksi Matematika} & \textbf{Loop Invariant} \\ \hline
Langkah basis: $P(n_0)$ benar & Inisialisasi: $I$ benar sebelum loop \\ \hline
Langkah induksi: $P(k) \Rightarrow P(k+1)$ & Pemeliharaan: $I$ sebelum iterasi $\Rightarrow$ $I$ setelah iterasi \\ \hline
Kesimpulan: $P(n)$ untuk semua $n$ & Terminasi: $I \wedge \neg B$ setelah loop berakhir \\ \hline
\end{tabular}
\end{center}

\subsection{Contoh: Penjumlahan Array}

\begin{contoh}
\textbf{Verifikasi Penjumlahan Elemen Array}

Program:
\begin{verbatim}
    sum := 0
    i := 0
    while i < n do
        sum := sum + A[i]
        i := i + 1
\end{verbatim}

\textbf{Spesifikasi}:
\begin{itemize}
    \item Prekondisi: $n \geq 0$ dan $A[0..n-1]$ terdefinisi
    \item Postkondisi: $sum = \sum_{j=0}^{n-1} A[j]$
\end{itemize}

\textbf{Loop Invariant}: $I \equiv \left(sum = \sum_{j=0}^{i-1} A[j]\right) \wedge (0 \leq i \leq n)$

\textbf{Verifikasi tiga syarat}:

\textbf{1. Inisialisasi}: Sebelum loop, $sum = 0$ dan $i = 0$. Konvensi \textit{penjumlahan kosong} (empty sum) mendefinisikan $\sum_{j=0}^{-1} A[j] = 0$. Maka $sum = 0$ dan $0 \leq 0 \leq n$. $\checkmark$

\textbf{2. Pemeliharaan}: Asumsikan $I$ benar dan $i < n$.
\begin{itemize}
    \item Setelah $sum := sum + A[i]$: $sum_{baru} = \sum_{j=0}^{i-1} A[j] + A[i] = \sum_{j=0}^{i} A[j]$
    \item Setelah $i := i + 1$: $i_{baru} = i + 1$
    \item Maka: $sum_{baru} = \sum_{j=0}^{i_{baru}-1} A[j]$ dan $0 \leq i_{baru} \leq n$. Invarian terjaga. $\checkmark$
\end{itemize}

\textbf{3. Terminasi}: Ketika loop berakhir, $\neg(i < n)$ yaitu $i \geq n$. Bersama invarian ($i \leq n$), diperoleh $i = n$. Maka $sum = \sum_{j=0}^{n-1} A[j]$, yang merupakan postkondisi. $\checkmark$ $\blacksquare$
\end{contoh}

\begin{catatan}
Menemukan loop invariant yang tepat seringkali menjadi bagian tersulit dari verifikasi program. Strategi yang berguna: invariant biasanya merupakan ``versi parsial'' dari postkondisi --- yaitu postkondisi yang berlaku untuk sebagian data yang sudah diproses \cite{wiki_loop_invariant}.
\end{catatan}
